% !PRINTED 183
\weberpage{183}

The representation is, as (2) shows, possible in infinitely many different ways, since the entirely arbitrary quantities $x'$ still occur in this formula.

A representation by fewer than $m$ squares is possible only when the determinant of the function $\varphi$ vanishes.

Since a square root occurs in the expression (2) for $Y_1$, $Y_1$ can be real or imaginary, even when the coefficients of $\varphi$ and the $x$ and $x'$ are assumed to be real.

But if, in the case where $Y_1$ is imaginary, we put $iY_1$ in place of $Y_1$, we can also state our theorem as follows:

\gesperrt{A real quadratic form in $m$ variables can be represented as a sum of at most $m$ positive or negative squares of real linear functions of the $x$}\footnote{For the general theory of the linear transformation of forms of the second degree, besides the investigations of \gesperrt{Jacobi} and \gesperrt{Hesse} (\gesperrt{Hesse}, ``Vorlesungen über analytische Geometrie des Raumes'', third edition, edited by \gesperrt{Gundelfinger}, Leipzig 1876), the works of \gesperrt{Weierstrass} and \gesperrt{Kronecker} (Monatsberichte der Berliner Akademie 1868, 1874) are of particular importance. A thorough historical account of the development of the question is given in the ``Bericht über den gegenwärtigen Stand der Invariantentheorie'' by \gesperrt{Franz Meyer} (first Jahresbericht der deutschen Mathematiker-Vereinigung, Berlin 1892).}.

\begin{center}
\S. 58.

\gesperrt{Law of inertia of quadratic forms.}
\end{center}

For the decomposition of a quadratic form into squares the following theorem now holds, which received from \gesperrt{Sylvester} the name of the law of inertia of quadratic forms (Philos. Mag. 1852).

\gesperrt{However one may decompose a real quadratic form $\varphi(x)$ into a sum of positive and negative squares of linear functions, the number of the positive and of the negative among these squares, and hence also their total number, is always the same, provided that no linear dependence exists among these linear functions.}

The proof of this important theorem is very simple.

Let
%CONT>>
